% !TEX program = pdflatex
% ==========================================================
% CLASE BEAMER MASTER
% Fundamentos de Matemática Básica - Agronomía
% Tema: Potencias
% Universidad Santo Tomás - Talca
% ==========================================================

\newif\ifhandout
%\handoutfalse
\handouttrue

\ifhandout
\documentclass[aspectratio=169,10pt,handout]{beamer}
\else
\documentclass[aspectratio=169,10pt]{beamer}
\fi

% ----------------------------------------------------------
% PAQUETES
% ----------------------------------------------------------
\usepackage[spanish,es-noshorthands]{babel}
\usepackage[T1]{fontenc}
\usepackage[utf8]{inputenc}
\usepackage{lmodern}
\usepackage{amsmath,amsfonts,amssymb}
\usepackage{ragged2e}
\usepackage{enumitem}
\usepackage{multicol}
\usepackage{array}
\usepackage{booktabs}
\usepackage{tikz}
\usetikzlibrary{positioning,calc}
\usepackage{fontawesome5}
\usepackage{colortbl}

% ----------------------------------------------------------
% COLORES
% ----------------------------------------------------------
\definecolor{USTBlue}{RGB}{0,70,127}
\definecolor{USTBlueDark}{RGB}{0,45,85}
\definecolor{USTTeal}{RGB}{0,153,117}
\definecolor{USTGray}{RGB}{120,120,120}
\definecolor{USTSoft}{RGB}{248,249,250}
\definecolor{USTText}{RGB}{35,40,45}
\definecolor{USTGold}{RGB}{196,154,70}
\definecolor{USTLightBlue}{RGB}{227,239,248}
\definecolor{USTLightTeal}{RGB}{226,245,240}
\definecolor{USTRedSoft}{RGB}{252,235,235}
\definecolor{USTLine}{RGB}{220,226,232}
\definecolor{USTLightGold}{RGB}{247,240,223}

% ----------------------------------------------------------
% TEMA
% ----------------------------------------------------------
\usetheme{Madrid}
\usecolortheme[named=USTBlue]{structure}
\setbeamertemplate{navigation symbols}{}

% ----------------------------------------------------------
% AJUSTES VISUALES
% ----------------------------------------------------------
\setlist[itemize]{itemsep=0.4em}
\setlist[enumerate]{itemsep=0.4em}

\setbeamercolor{normal text}{fg=USTText,bg=white}
\setbeamercolor{title}{fg=USTBlueDark}
\setbeamercolor{subtitle}{fg=USTTeal}
\setbeamercolor{frametitle}{fg=white,bg=USTBlueDark}

\setbeamercolor{block title}{bg=USTLightBlue,fg=USTBlueDark}
\setbeamercolor{block body}{bg=USTSoft,fg=black}

\setbeamercolor{block title example}{bg=USTLightTeal,fg=USTBlueDark}
\setbeamercolor{block body example}{bg=USTSoft,fg=black}

\setbeamercolor{block title alerted}{bg=USTRedSoft,fg=black}
\setbeamercolor{block body alerted}{bg=red!2,fg=black}

% ----------------------------------------------------------
% FRAMETITLE
% ----------------------------------------------------------
\setbeamertemplate{frametitle}{
	\vspace*{-0.03cm}
	\begin{beamercolorbox}[wd=\paperwidth,ht=0.92cm,dp=0.18cm,leftskip=0.45cm,rightskip=0.25cm]{frametitle}
		{\large\bfseries \insertframetitle}
	\end{beamercolorbox}
}

% ----------------------------------------------------------
% PIE DE PÁGINA
% ----------------------------------------------------------
\setbeamercolor{author in head/foot}{bg=USTSoft,fg=USTGray}
\setbeamercolor{date in head/foot}{bg=USTSoft,fg=USTGray}
\setbeamercolor{title in head/foot}{bg=USTSoft,fg=USTGray}

\setbeamertemplate{footline}{%
	\leavevmode%
	\hbox{%
		\begin{beamercolorbox}[wd=.82\paperwidth,ht=2.8ex,dp=1ex,leftskip=0.4cm]{author in head/foot}%
			\scriptsize
			\textbf{Fundamentos de Matemática Básica}
			\quad|\quad
			Agronomía
			\quad|\quad
			Universidad Santo Tomás Talca
		\end{beamercolorbox}%
		\begin{beamercolorbox}[wd=.18\paperwidth,ht=2.8ex,dp=1ex,rightskip=0.35cm plus1fil]{date in head/foot}%
			\scriptsize \insertframenumber/\inserttotalframenumber
		\end{beamercolorbox}%
	}%
}

% ----------------------------------------------------------
% COMANDOS
% ----------------------------------------------------------
\newcommand{\destacar}[1]{\textcolor{USTBlueDark}{\textbf{#1}}}
\newcommand{\alerta}[1]{\textcolor{red!70!black}{\textbf{#1}}}
\newcommand{\pp}{\pause}

\newcommand{\iconoInicio}{\faPlayCircle}
\newcommand{\iconoDesarrollo}{\faBrain}
\newcommand{\iconoPractica}{\faPenSquare}
\newcommand{\iconoAgro}{\faSeedling}
\newcommand{\iconoCierre}{\faCheckCircle}
\newcommand{\iconoIdea}{\faLightbulb}
\newcommand{\iconoError}{\faExclamationTriangle}
\newcommand{\iconoTiempo}{\faClock}
\newcommand{\iconoMeta}{\faBullseye}
\newcommand{\iconoEval}{\faClipboardCheck}
\newcommand{\iconoTarea}{\faBook}
\newcommand{\iconoSalida}{\faDoorOpen}

\newenvironment{metaBox}
{\begin{block}{\iconoMeta\ \ Meta de aprendizaje}}
	{\end{block}}

\newenvironment{ideaBox}
{\begin{exampleblock}{\iconoIdea\ \ Idea clave}}
	{\end{exampleblock}}

% ----------------------------------------------------------
% PORTADA MASTER
% ----------------------------------------------------------
\newcommand{\PortadaProfesionalUST}{%
	\begin{frame}[plain]
		\begin{tikzpicture}[remember picture,overlay]
			\fill[white] (current page.south west) rectangle (current page.north east);
			\fill[USTBlueDark] (current page.north west) rectangle ([yshift=-1.45cm]current page.north east);
			\fill[USTTeal] ([yshift=-1.45cm]current page.north west) rectangle ([yshift=-1.62cm]current page.north east);
			\fill[USTSoft] (current page.south west) rectangle ([xshift=2.25cm]current page.north west);
			
			\fill[USTBlue!7] ([xshift=1.10cm,yshift=-2.0cm]current page.north west) circle (0.78cm);
			\fill[USTTeal!10] ([xshift=1.70cm,yshift=-4.2cm]current page.north west) circle (0.45cm);
			\fill[USTGold!14] ([xshift=-1.2cm,yshift=1cm]current page.south east) circle (1.08cm);
			
			\fill[black!7,rounded corners=8pt]
			([xshift=-5.1cm,yshift=-3.15cm]current page.center) rectangle
			([xshift=5.0cm,yshift=2.75cm]current page.center);
			
			\fill[white,rounded corners=8pt]
			([xshift=-5.25cm,yshift=-3.00cm]current page.center) rectangle
			([xshift=4.85cm,yshift=2.90cm]current page.center);
			
			\fill[USTTeal] ([xshift=-5.25cm,yshift=-3.00cm]current page.center) rectangle
			([xshift=-4.95cm,yshift=2.90cm]current page.center);
			
			\node[anchor=north east,text=white]
			at ([xshift=-0.7cm,yshift=-0.52cm]current page.north east)
			{\small\bfseries Universidad Santo Tomás \quad Ciencias Básicas Talca};
			
			\node[align=center] at ([xshift=0.10cm,yshift=-0.05cm]current page.center){%
				\begin{minipage}{0.68\paperwidth}
					\centering
					{\Large\bfseries\color{USTText}\insertauthor\par}
					\vspace{0.75cm}
					{\fontsize{24}{28}\selectfont\bfseries\color{USTBlueDark}\inserttitle\par}
					\vspace{0.28cm}
					{\large\bfseries\color{USTTeal}\insertsubtitle\par}
					\vspace{0.58cm}
					{\normalsize\color{USTGray}\faSeedling\ Agronomía \quad \faUniversity\ Universidad Santo Tomás\par}
					\vspace{0.65cm}
					{\small\color{USTGray}\insertdate\par}
				\end{minipage}
			};
			
			\draw[USTGray!35,line width=0.35pt]
			([xshift=0.9cm,yshift=0.72cm]current page.south west) --
			([xshift=-0.9cm,yshift=0.72cm]current page.south east);
			
			\node[anchor=south,text=USTGray]
			at ([yshift=0.2cm]current page.south)
			{\scriptsize
				\textbf{Fundamentos de Matemática Básica}
				\;|\;
				Clase MASTER
				\;|\;
				Potencias};
		\end{tikzpicture}
	\end{frame}
}

% ----------------------------------------------------------
% FRAME DE SECCIÓN
% ----------------------------------------------------------
\newcommand{\SectionFrame}[3]{%
	\begin{frame}[plain]
		\begin{tikzpicture}[remember picture,overlay]
			\fill[USTBlueDark] (current page.south west) rectangle (current page.north east);
			\fill[USTTeal] (current page.south west) rectangle ([yshift=0.23cm]current page.north east);
			\fill[white,opacity=0.08] ([xshift=-1cm,yshift=1cm]current page.south east) circle (1.8cm);
			\fill[white,opacity=0.06] ([xshift=1.3cm,yshift=-0.8cm]current page.north west) circle (1.25cm);
			
			\node[anchor=center,text=white] at ([yshift=0.9cm]current page.center)
			{{\fontsize{34}{36}\selectfont #1}};
			
			\node[anchor=center,text=white] at (current page.center)
			{{\LARGE\bfseries #2}};
			
			\node[anchor=center,text=white!90] at ([yshift=-1.0cm]current page.center)
			{{\large #3}};
		\end{tikzpicture}
	\end{frame}
}

% ----------------------------------------------------------
% DATOS
% ----------------------------------------------------------
\title{Fundamentos de Matemática Básica}
\subtitle{Clase MASTER: Potencias}
\author{Prof. Luis González Alcaino}
\institute{Universidad Santo Tomás}
\date{Año académico 2026}

% ==========================================================
% DOCUMENTO
% ==========================================================
\begin{document}
	
	\PortadaProfesionalUST
	
	% ----------------------------------------------------------
	\begin{frame}{Panorama general}
		\begin{metaBox}
			Aplicar el concepto de potencia y sus propiedades en ejercicios matemáticos y situaciones contextualizadas en Agronomía.
		\end{metaBox}
		
		\pp
		
		\begin{block}{\faListOl\ \ Ruta de la sesión}
			\begin{enumerate}
				\item Actividad diagnóstica
				\item Desarrollo conceptual
				\item Evaluación formativa breve
				\item Aplicación agronómica
				\item Práctica autónoma
				\item Ticket de salida y tarea
			\end{enumerate}
		\end{block}
	\end{frame}
	
	% ----------------------------------------------------------
	\begin{frame}{\iconoTiempo\ \ Distribución del tiempo}
		\centering
		\begin{tabular}{>{\raggedright\arraybackslash}p{4.2cm} >{\centering\arraybackslash}p{2cm} >{\raggedright\arraybackslash}p{4.8cm}}
			\toprule
			\textbf{Momento} & \textbf{Tiempo} & \textbf{Propósito} \\
			\midrule
			Diagnóstico inicial & 10 min & Activar saberes previos \\
			Desarrollo conceptual & 20 min & Comprender definición y propiedades \\
			Práctica guiada & 15 min & Resolver con apoyo docente \\
			Evaluación formativa & 10 min & Verificar comprensión inmediata \\
			Aplicación agronómica & 10 min & Conectar con el contexto profesional \\
			Práctica autónoma & 10 min & Consolidar individualmente \\
			Cierre y ticket de salida & 5 min & Sintetizar y reflexionar \\
			\bottomrule
		\end{tabular}
	\end{frame}
	
	% ----------------------------------------------------------
	\SectionFrame{\iconoInicio}{Inicio}{Actividad diagnóstica y activación}
	
	% ----------------------------------------------------------
	\begin{frame}{\iconoInicio\ \ Actividad diagnóstica}
		\begin{block}{Responde mentalmente o en tu cuaderno}
			\begin{enumerate}
				\item ¿Cómo escribirías \(3\cdot3\cdot3\cdot3\) usando potencia?
				\item ¿Qué significa el exponente en \(5^2\)?
				\item ¿Cuánto vale \(10^3\)?
				\item ¿Qué fenómeno agrícola podría crecer de manera repetida?
			\end{enumerate}
		\end{block}
	\end{frame}
	
	% ----------------------------------------------------------
	\begin{frame}{\iconoInicio\ \ Puesta en común}
		\begin{exampleblock}{Posibles respuestas}
			\begin{itemize}
				\item \(3\cdot3\cdot3\cdot3 = 3^4\)
				\item En \(5^2\), el exponente indica que el 5 se repite dos veces
				\item \(10^3=1000\)
				\item Ejemplos: bacterias, semillas, unidades experimentales
			\end{itemize}
		\end{exampleblock}
		
		\pp
		
		\begin{ideaBox}
			Las potencias permiten expresar repetición, crecimiento y escala.
		\end{ideaBox}
	\end{frame}
	
	% ----------------------------------------------------------
	\SectionFrame{\iconoDesarrollo}{Desarrollo}{Concepto y propiedades}
	
	% ----------------------------------------------------------
	\begin{frame}{\iconoDesarrollo\ \ Concepto de potencia}
		\begin{block}{Definición}
			\[
			a^n=\underbrace{a\cdot a\cdot a\cdots a}_{n\text{ veces}}
			\]
		\end{block}
		
		\pp
		
		\begin{columns}[T,totalwidth=\textwidth]
			\begin{column}{0.47\textwidth}
				\begin{block}{Base}
					Es el número que se repite.
				\end{block}
			\end{column}
			\begin{column}{0.47\textwidth}
				\begin{exampleblock}{Exponente}
					Indica cuántas veces se repite la base.
				\end{exampleblock}
			\end{column}
		\end{columns}
		
		\pp
		
		\begin{exampleblock}{Ejemplo}
			\[
			5^3=5\cdot5\cdot5=125
			\]
		\end{exampleblock}
	\end{frame}
	
	% ----------------------------------------------------------
	\begin{frame}{\iconoDesarrollo\ \ Propiedades fundamentales I}
		\begin{block}{Producto de igual base}
			\[
			a^m\cdot a^n=a^{m+n}
			\]
		\end{block}
		\begin{exampleblock}{Ejemplo}
			\[
			2^3\cdot2^4=2^7=128
			\]
		\end{exampleblock}
		
		\pp
		
		\begin{block}{Cociente de igual base}
			\[
			\frac{a^m}{a^n}=a^{m-n},\qquad a\neq0
			\]
		\end{block}
		\begin{exampleblock}{Ejemplo}
			\[
			\frac{5^6}{5^2}=5^4=625
			\]
		\end{exampleblock}
	\end{frame}
	
	% ----------------------------------------------------------
	\begin{frame}{\iconoDesarrollo\ \ Propiedades fundamentales II}
		\begin{block}{Potencia de una potencia}
			\[
			(a^m)^n=a^{m\cdot n}
			\]
		\end{block}
		\begin{exampleblock}{Ejemplo}
			\[
			(3^2)^4=3^8=6561
			\]
		\end{exampleblock}
		
		\pp
		
		\begin{block}{Potencia de un producto}
			\[
			(ab)^n=a^n b^n
			\]
		\end{block}
		\begin{exampleblock}{Ejemplo}
			\[
			(2\cdot5)^3=2^3\cdot5^3=1000
			\]
		\end{exampleblock}
	\end{frame}
	
	% ----------------------------------------------------------
	\begin{frame}{\iconoDesarrollo\ \ Propiedades fundamentales III}
		\begin{block}{Potencia de un cociente}
			\[
			\left(\frac{a}{b}\right)^n=\frac{a^n}{b^n},\qquad b\neq0
			\]
		\end{block}
		\begin{exampleblock}{Ejemplo}
			\[
			\left(\frac{2}{3}\right)^3=\frac{8}{27}
			\]
		\end{exampleblock}
		
		\pp
		
		\begin{block}{Exponente cero y negativo}
			\[
			a^0=1,\ a\neq0
			\qquad
			a^{-n}=\frac{1}{a^n}
			\]
		\end{block}
		\begin{exampleblock}{Ejemplos}
			\[
			7^0=1
			\qquad
			2^{-3}=\frac{1}{8}
			\]
		\end{exampleblock}
	\end{frame}
	
	% ----------------------------------------------------------
	\begin{frame}{\iconoError\ \ Errores frecuentes}
		\begin{columns}[T,totalwidth=\textwidth]
			\begin{column}{0.48\textwidth}
				\begin{alertblock}{Error 1}
					\[
					a^m+a^n \neq a^{m+n}
					\]
					\[
					2^2+2^3=12
					\]
					\[
					2^5=32
					\]
				\end{alertblock}
			\end{column}
			\begin{column}{0.48\textwidth}
				\begin{alertblock}{Error 2}
					\[
					(a+b)^n \neq a^n+b^n
					\]
					\[
					(2+3)^2=25
					\]
					\[
					2^2+3^2=13
					\]
				\end{alertblock}
			\end{column}
		\end{columns}
	\end{frame}
	
	% ----------------------------------------------------------
	\SectionFrame{\iconoPractica}{Práctica guiada}{Resolución paso a paso}
	
	% ----------------------------------------------------------
	\begin{frame}{\iconoPractica\ \ Ejercicio guiado 1}
		\begin{block}{Problema}
			Simplificar:
			\[
			2^4\cdot2^3
			\]
		\end{block}
		
		\only<2>{
			\begin{exampleblock}{Paso 1}
				Misma base:
				\[
				2^4\cdot2^3=2^{4+3}
				\]
			\end{exampleblock}
		}
		
		\only<3>{
			\begin{exampleblock}{Paso 2}
				\[
				2^7=128
				\]
			\end{exampleblock}
		}
	\end{frame}
	
	% ----------------------------------------------------------
	\begin{frame}{\iconoPractica\ \ Ejercicio guiado 2}
		\begin{block}{Problema}
			Resolver:
			\[
			(3^2)^3
			\]
		\end{block}
		
		\only<2>{
			\begin{exampleblock}{Paso 1}
				Multiplicamos exponentes:
				\[
				(3^2)^3=3^{2\cdot3}
				\]
			\end{exampleblock}
		}
		
		\only<3>{
			\begin{exampleblock}{Paso 2}
				\[
				3^6=729
				\]
			\end{exampleblock}
		}
	\end{frame}
	
	% ----------------------------------------------------------
	\SectionFrame{\iconoEval}{Evaluación formativa}{Verificación rápida de comprensión}
	
	% ----------------------------------------------------------
	\begin{frame}{\iconoEval\ \ Chequeo rápido}
		Responde individualmente:
		\begin{enumerate}
			\item \(4^2\cdot4^3=\ ?\)
			\item \(\dfrac{6^5}{6^2}=\ ?\)
			\item \((2^3)^2=\ ?\)
			\item \(5^0=\ ?\)
		\end{enumerate}
	\end{frame}
	
	% ----------------------------------------------------------
	\begin{frame}{\iconoEval\ \ Retroalimentación inmediata}
		\begin{exampleblock}{Respuestas}
			\begin{enumerate}
				\item \(4^2\cdot4^3=4^5\)
				\item \(\dfrac{6^5}{6^2}=6^3\)
				\item \((2^3)^2=2^6\)
				\item \(5^0=1\)
			\end{enumerate}
		\end{exampleblock}
		
		\begin{ideaBox}
			Si hay seguridad en estas cuatro respuestas, la base conceptual está bien encaminada.
		\end{ideaBox}
	\end{frame}
	
	% ----------------------------------------------------------
	\SectionFrame{\iconoAgro}{Aplicación agronómica}{Conexión con el contexto profesional}
	
	% ----------------------------------------------------------
	\begin{frame}{\iconoAgro\ \ Caso 1: semillas por superficie}
		\begin{block}{Situación}
			En una parcela experimental se usan:
			\[
			2^5
			\]
			semillas por metro cuadrado.
			
			La superficie cultivada es:
			\[
			2^3
			\]
			metros cuadrados.
		\end{block}
		
		\pp
		
		\begin{exampleblock}{Resolución}
			\[
			2^5\cdot2^3=2^8=256
			\]
			Se utilizan \textbf{256 semillas}.
		\end{exampleblock}
	\end{frame}
	
	% ----------------------------------------------------------
	\begin{frame}{\iconoAgro\ \ Caso 2: fertilización}
		\begin{block}{Situación}
			Se aplican:
			\[
			5\times10^2 \text{ gramos por hectárea}
			\]
			en una superficie de:
			\[
			10^3 \text{ hectáreas}
			\]
		\end{block}
		
		\pp
		
		\begin{exampleblock}{Resolución}
			\[
			(5\times10^2)(10^3)=5\times10^5=500000
			\]
			Se requieren \textbf{500000 gramos}.
		\end{exampleblock}
	\end{frame}
	
	% ----------------------------------------------------------
	\begin{frame}{\iconoAgro\ \ Caso 3: concentración de micronutrientes}
		\begin{block}{Situación}
			Una solución contiene:
			\[
			3\times10^{-2}\ \text{g}
			\]
			de micronutriente por litro.
			
			Si se preparan:
			\[
			10^3
			\]
			litros, ¿cuántos gramos se requieren?
		\end{block}
		
		\pp
		
		\begin{exampleblock}{Resolución}
			\[
			(3\times10^{-2})(10^3)=3\times10^1=30
			\]
			Se necesitan \textbf{30 gramos}.
		\end{exampleblock}
	\end{frame}
	
	% ----------------------------------------------------------
	\SectionFrame{\iconoPractica}{Práctica autónoma}{Consolidación individual}
	
	% ----------------------------------------------------------
	\begin{frame}{\iconoPractica\ \ Ejercicios propuestos}
		\begin{enumerate}
			\item \(2^5\cdot2^2\)
			\item \(\dfrac{3^7}{3^4}\)
			\item \((4^2)^3\)
			\item \((2\cdot3)^2\)
			\item \(\left(\dfrac{3}{5}\right)^2\)
			\item \(8^0\)
			\item \(2^{-4}\)
			\item En un cultivo hay \(3^4\) plantas por sector y \(3^2\) sectores. ¿Cuántas plantas hay en total?
		\end{enumerate}
	\end{frame}
	
	% ----------------------------------------------------------
	\begin{frame}{\iconoPractica\ \ Respuestas}
		\begin{multicols}{2}
			\begin{enumerate}
				\item \(2^5\cdot2^2=2^7=128\)
				\item \(\dfrac{3^7}{3^4}=3^3=27\)
				\item \((4^2)^3=4^6=4096\)
				\item \((2\cdot3)^2=36\)
				\item \(\left(\dfrac{3}{5}\right)^2=\dfrac{9}{25}\)
				\item \(8^0=1\)
				\item \(2^{-4}=\dfrac{1}{16}\)
				\item \(3^4\cdot3^2=3^6=729\)
			\end{enumerate}
		\end{multicols}
	\end{frame}
	
	% ----------------------------------------------------------
	\SectionFrame{\iconoCierre}{Cierre}{Síntesis, ticket de salida y proyección}
	
	% ----------------------------------------------------------
	\begin{frame}{\iconoCierre\ \ Síntesis}
		\begin{block}{Ideas clave}
			\begin{itemize}
				\item Una potencia representa una multiplicación repetida.
				\item La base se repite; el exponente indica cuántas veces.
				\item Las propiedades simplifican cálculos.
				\item Las potencias son útiles en contextos agrícolas reales.
			\end{itemize}
		\end{block}
	\end{frame}
	
	% ----------------------------------------------------------
	\begin{frame}{\iconoSalida\ \ Ticket de salida}
		\begin{exampleblock}{Responder antes de terminar la clase}
			\begin{enumerate}
				\item ¿Qué propiedad usarías para resolver \(7^2\cdot7^5\)?
				\item ¿Cuánto vale \(10^{-2}\)?
				\item Menciona una aplicación de las potencias en Agronomía.
			\end{enumerate}
		\end{exampleblock}
	\end{frame}
	
	% ----------------------------------------------------------
	\begin{frame}{\iconoEval\ \ Rúbrica breve de desempeño}
		\centering
		\begin{tabular}{>{\raggedright\arraybackslash}p{5.2cm} >{\centering\arraybackslash}p{2.2cm} >{\centering\arraybackslash}p{2.2cm}}
			\toprule
			\textbf{Criterio} & \textbf{Logrado} & \textbf{En proceso} \\
			\midrule
			Reconoce base y exponente & Sí & Parcial \\
			Aplica propiedades correctamente & Sí & Parcial \\
			Evita errores frecuentes & Sí & Parcial \\
			Interpreta una situación agronómica & Sí & Parcial \\
			\bottomrule
		\end{tabular}
	\end{frame}
	
	% ----------------------------------------------------------
	\begin{frame}{\iconoTarea\ \ Tarea sugerida}
		\begin{block}{Para la próxima clase}
			Resolver en el cuaderno:
			\begin{enumerate}
				\item \(5^3\cdot5^4\)
				\item \(\dfrac{10^7}{10^3}\)
				\item \((2^4)^2\)
				\item \(\left(\dfrac{4}{7}\right)^2\)
				\item Crear un ejemplo propio de uso de potencias en Agronomía.
			\end{enumerate}
		\end{block}
		
		\begin{ideaBox}
			La próxima sesión puede enlazarse naturalmente con notación científica y operaciones con potencias de 10.
		\end{ideaBox}
	\end{frame}
	
	% ----------------------------------------------------------
	\begin{frame}{\iconoCierre\ \ Reflexión final}
		\begin{exampleblock}{Pregunta final}
			¿Por qué comprender potencias puede ayudarte a interpretar mejor datos, escalas y fenómenos de crecimiento en el ámbito agrícola?
		\end{exampleblock}
		
		\vspace{0.5cm}
		\begin{center}
			{\Large\color{USTTeal}\faSeedling}
			
			\vspace{0.25cm}
			{\normalsize\textcolor{USTGray}{La matemática básica no solo sirve para calcular: sirve para comprender, modelar y decidir.}}
		\end{center}
	\end{frame}
	
\end{document}